\section{Chain Complexes}

\subsection{Complexes of $R$-modules}

\bdefn

    Let $R$ be a ring.
    A {\emphcolor chain complex} of $R$ modules is a sequence of $R$-modules $C_\bul=\set{C_n}_{n\in\bZ}$ along with maps $d_n\colon C_n\to C_{n-1}$
    such that $d_{n-1}\circ d_n\colon C_n\to C_{n-2}$ is zero for all $n$.
    These maps $d_n$ are called the {\emphcolor differentials} of $C_\bul$.

\edefn

We often omit the subscripts for the differentials and the chain complex itself; we write $d$ and $C$ instead of $d_n$ and $C_\bul$ respectively.

Note that $d_n\circ d_{n+1}=0$ means that $\im d_{n+1}\subseteq\ker d_n$.
This justifies the following definition:

\bdefn

    Given a chain complex $C$, define
    \benum
        \item the kernel of $d_n\colon C_n\to C_{n-1}$ to be $Z_n=Z_n(C)$, the {\emphcolor $n$-cycles} of $C$;
        \item the image of $d_{n+1}\colon C_{n+1}\to C_n$ to be $B_n=B_n(C)$, the {\emphcolor $n$-boundaries} of $C$;
        \item the {\emphcolor $n$th homology module} of $C$ to be their quotient: $Z_n/B_n$.
    \eenum

\edefn

\bexam

    Let $C_n=\bZ/8\bZ$ for $n\geq0$ and $C_n=0$ for $n<0$.
    Define $d_n$ to map $x$ to $4x$ (modulo $8$).

    Then we have $d\circ d$ maps $x$ to $8x\equiv0$, so this forms a chain complex (of $\bZ/8\bZ$-modules, say).
    And $Z_n$ is $\bZ/2\bZ$, $B_n$ is $\bZ/4\bZ$.
    Thus we have $H_n$ is (isomorphic to) $\bZ/2\bZ$.

\eexam

We would like to build a category of chain complexes of $R$-modules.
To do so we must define morphisms of chain complexes; there is a clear way to do so.

\bdefn

    A {\emphcolor morphism} of chain complexes $u\colon C\to D$ is a collection of $R$-module morphisms $u_n\colon C_n\to D_n$
    which commute with the differential maps: $u_{n-1}\circ d_n=d_{n-1}\circ u_n$.

    The category $\Ch(\Mod_R)$ is the category of chain complexes of (right) $R$-modules.

\edefn

\bprop

    A morphism of chain complexes $u\colon C\to D$ maps boundaries to boundaries and cycles to cycles.
    Thus defines maps $H_n(C)\to H_n(D)$.
    Furthermore, for each $n\in\bZ$, $H_n$ defines a functor $\Ch(\Mod_R)\to\Mod_R$.

\eprop

\bproof

    We have that $udx=dux$, thus for $dx\in B(C)$ its image under $u$ is in $B(D)$.
    Similarly if $dx=0$, then $udx=dux=0$, in particular $ux\in Z(D)$.
    And so $u_n$ factors through $B_n(C)$ and restricts to cycles, giving us maps $u_n\colon H_n(C)\to H_n(D)$.

    This is a functor, as quotienting a morphism is.
    \qed

\eproof

\bexam

    Let $k$ be a field, and consider $k$-vector spaces $\set{B_n,H_n}_{n\in\bZ}$.
    Define $C_n=B_n\oplus H_n\oplus B_{n-1}$ and define $d_n\colon C_n\to C_{n-1}$ by projection onto $B_{n-1}\subseteq C_{n-1}$.
    This is a chain complex: $d_n$ projects onto $B_{n-1}$ which is orthogonal to $B_{n-2}$, onto which $d_{n-1}$ projects.
    Further note that $Z_n(C)=B_n\oplus H_n$ and $B_n(C)=B_n$, thus $H_n(C)=H_n$; justifying the choice of names.

    If $V$ is a chain complex of $k$-vector spaces, then note that $V_n\cong B_n(V)\oplus H_n(V)\oplus B_{n-1}(V)$.
    Thus all chain complexes of $k$-vector spaces can be constructed in this manner.

\eexam

\bexam

    Let $C$ be a chain complex of $R$-modules, and $A$ an $R$-module.
    Then we define $d'_n\colon\hom_R(A,C_n)\to\hom_R(A,C_{n-1})$ by mapping $f$ to $d_nf$.
    This forms a chain complex clearly.

    Now if $A=Z_n(C)$ for some $n$, suppose $H_n(\hom_R(Z_n,C))=0$.
    This means that $B_n(\hom_R(Z_n,C))=Z_n(\hom_R(Z_n,C))$, i.e.\ if $df=0$ then $f=dg$ for all $f\colon Z_n\to C_n$.
    In particular for inclusion $\iota\colon Z_n\to C_n$; $d\iota=0$ so $\iota=df$ for some $f\colon Z_n\to C_{n+1}$.
    But this means that $Z_n\subseteq B_n$ (and thus there is equality): $Z_n=\im\iota=\im df\subseteq\im d=B_n$, and so $H_n(C)=0$.

\eexam

\bdefn

    A morphism of chain complexes $C\to D$ is a {\emphcolor quasi-isomorphism} if the maps $H_n(C)\to H_n(D)$ are all isomorphisms.

\edefn

There is a clear chain complex for all $R$: the zero chain complex, consisting of only the zero $R$-module.
This complex is a zero object in $\Ch(\Mod_R)$: it is both initial and final.

\bprop

    The following are equivalent for every chain complex $C$:
    \benum
        \item $C$ is exact;
        \item $C$ is {\emphcolor acyclic}: $H_n(C)=0$ for all $n$;
        \item The map $0\to C$ is a quasi-isomorphism.
    \eenum

\eprop

\bproof

    $(1)$ and $(2)$ are clearly equivalent: $H_n(C)=0$ iff $B_n(C)=Z_n(C)$, which is just what it means for $C$ to be exact.
    These are also clearly equivalent to $(3)$, by definition.
    \qed

\eproof

\bdefn

    A {\emphcolor cochain complex} of $R$-modules is $C^\bul$, consisting of a collection of $R$-modules $\set{C^n}_{n\in\bZ}$ along with differential maps
    $d^n\colon C^n\to C^{n+1}$ such that $d\circ d=0$.
    As before, $Z^n(C^\bul)=\ker d^n$ is the module of {\emphcolor $n$-cocycles} and $B^n(C^\bul)=\im d^{n-1}\subseteq C^n$ is the module of {\emphcolor $n$-coboundaries}.
    $H^n(C^\bul)=Z^n(C^\bul)/B^n(C^\bul)$ is the {\emphcolor $n$th cohomology module} of $C^\bul$.

    Morphisms of cochain complexes are defined the same.

\edefn

If $C_\bul$ is a chain complex we can define a cochain complex $C^\bul$ simply by $C^n=C_{-n}$ and $d^n=d_{-n}$.

\bdefn

    A chain complex $C$ is {\emphcolor bounded} if all but finitely many $C_n$ are zero.
    If $C_n=0$ for all $n$ other than $a\leq n\leq b$, then $C$ has {\emphcolor amplitude in $[a,b]$}.
    $C$ is {\emphcolor bounded above} (resp.\ {\emphcolor bounded below}) if there is a $b$ (resp.\ $a$) such that $C_n=0$ for $n>b$ (resp.\ $n<a$).

    The full subcategories of bounded, bounded above, bounded below chain complexes are denoted $\Ch_b,\Ch_-,\Ch_+$ respectively.
    The subcategory of non-negative chain complexes is denoted by $\Ch_{\geq0}$ (those chain complexes bounded below by $0$).

    Similar for cochain complexes.

\edefn

\subsection{Operations on chain complexes}

The purpose of this section is to show that chain complexes form an {\emphcolor Abelian category}.
Recall the following category-theoretic definitions:

\bdefn

    A {\emphcolor preadditive category} is a category whose hom-sets are given the structure of Abelian groups, such that composition distributes over addition: $h(g+g')f=hgf+hg'f$.
    An {\emphcolor additive functor} is a functor between two preadditive categories whose restrictions to the hom-sets forms an Abelian group homomorphism.

    Finite products and coproducts in a preadditive category necessarily coincide; they are called {\emphcolor biproducts}.
    An {\emphcolor additive category} is a preadditive category with all finite biproducts (equivalently a zero object and binary biproducts).

\edefn

$\Ch=\Ch(\Mod_R)$ is an additive category.
We can add chain complex morphisms in a clear manner: $(u+v)_n=u_n+v_n$.
Furthermore as discussed, $\Ch$ has a zero object.
$\Ch$ has all arbitrary direct products and coproducts (and thus all finite biproducts, so it is an additive category).

Indeed, let $\set{C_\alpha}_{\alpha\in I}$ be a collection of chain complexes of $R$-modules.
We define their product $\prod_\alpha C_\alpha$ to have as components the direct product of the components: the $n$th component is $\prod_\alpha C_{\alpha,n}$.
Similarly their coproduct $\bigoplus_\alpha C_\alpha$'s $n$th component is $\bigoplus_\alpha C_{\alpha,n}$.
The differential maps are defined coordinate-wise as well.
These satisfy the universal properties of products and coproducts respectively.

\bprop

    Products and coproducts commute with the homology functors.
    That is to say, $H_n(\prod_\alpha C_\alpha)\cong\prod_\alpha H_n(C_\alpha)$ and similar for coproducts.

\eprop

\bdefn

    A chain complex $B$ is a {\emphcolor subcomplex} of another chain complex $C$ if $B_n$ is a submodule of $C_n$ for every $n$, and $B$'s differential maps are the restrictions
    of $C$'s.

    If $B$ is a subcomplex of $C$, then we can define their {\emphcolor quotient complex} $C/B$ whose components are $C_n/B_n$ and whose differential maps are the factorizations of $C$ and $B$'s.

\edefn

In particular, if $f\colon B\to C$ is a chain map, the kernels $\set{\ker f_n}_n$ form a subcomplex of $B$ denoted $\ker f$.
And the cokernels $\set{\coker f_n}_n$ form a quotient complex of $C$ denoted $\coker f$.

\bdefn

    In a category with a zero object (more generally with zero morphisms), the {\emphcolor kernel} of a morphism $f$ is its equalizer with $0$.
    Its {\emphcolor cokernel} is its coequalizer with $0$.

    In a general category, a {\emphcolor monomorphism} is a morphism which is left-cancelative: $fg=fh$ implies $g=h$.
    An {\emphcolor epimorphism} is one which is right-cancelative: $gf=hf$ implies $g=h$.
    In a preadditive category monomorphisms are precisely those morphisms for which $fg=0$ implies $g=0$; epimorphisms are those for which $gf=0$ implies $g=0$.

\edefn

Our notion of kernels and cokernels in $\Ch$ coincide with the category-theoretic definitions.
Furthermore a monomorphism in $\Ch$ is precisely a chain morphism whose components are all monomorphisms, similar for epimorphisms.

\bdefn

    An {\emphcolor Abelian category} is an additive category such that
    \benum
        \item every map has a kernel and cokernel,
        \item every monomorphism is the kernel of its cokernel,
        \item every epimorphism is the cokernel of its kernel.
    \eenum

\edefn

The prototypical example of an Abelian category is $\Mod_R$ (in particular $\Ab=\Mod_\bZ$).

\bdefn

    Let $\C$ be a category, and $A\in\C$ an object.
    We say that two monomorphisms $u\colon X\to A$ and $v\colon Y\to A$ are equivalent if there exists an isomorphism $\phi\colon X\to Y$ such that $u=v\phi$.
    The equivalence classes of these monomorphisms are called {\emphcolor subobjects} of $A$.

    Similarly two epimorphisms $u\colon A\to X$ and $v\colon A\to Y$ are equivalent if there exists an isomorphism $\phi\colon X\to Y$ such that $v=\phi u$.
    The equivalence classes of these epimorphisms are called {\emphcolor quotient objects} of $A$.

\edefn

\bdefn

    In an Abelian category, the {\emphcolor image} of a morphism $f\colon B\to C$ is $\im f=\ker(\coker f)$, considered as a subobject of $C$.
    Its {\emphcolor coimage} is $\coim f=\coker(\ker f)$, viewed as a subobject of $B$.

\edefn

In $\Mod_R$ we define the image of $f\colon B\to C$ to be $\im f=\set{fb}[b\in B]$ (the coimage is $B/\ker f$).
This coincides with the category-theoretic definition, taking the monomorphism to be the inclusion $\im f\to C$.
Every map $f$ of $R$-modules factors as
$$ B\xto{\ e\ }\im f\xto{\ m\ }C $$
with $e$ epi and $m$ monic.

\bdefn

    In an Abelian category, a sequence of morphisms $A\xto{\ f\ }B\xto{\ g\ }C$ is {\emphcolor exact} at $B$ if $\im f=\ker g$ (as subobjects of $B$).

\edefn

Note that if $\A$ is an Abelian category, we can repeat the discussion of the previous section to define the category of chain complexes on $\A$, $\Ch(\A)$.
This is an additive category, and homology are additive functors $\Ch(\A)\to\A$.

\bthrm

    For an Abelian category $\A$, $\Ch(\A)$ is itself Abelian as well.

\ethrm

\bproof

    $\Ch(A)$ has binary biproducts and a zero object; their constructions are the same as before.
    Now if $f\colon B\to C$ is a chain map, then $f$ is monic iff $f_n\colon B_n\to C_n$ is monic for all $n$.
    Indeed, the composite $\ker f\to C$ (inclusion composed with $f$) is zero, and so if $f$ is monic then $\ker f=0$ (since $f\iota=0$ implies $\iota=0$), and conversely (recall $\ker f=0$ iff $f$ is monic).
    Thus if $f$ is monic, then $f$ is isomorphic to the kernel of $C\to C/B$.
    Similar for epimorphisms.
    \qed

\eproof

Since $\Ch(\A)$ is itself Abelian, we can form the category of chain complexes over $\Ch(\A)$.
These can be characterized more elegantly as follows:

\bdefn

    A {\emphcolor bicomplex} in $\A$ is a family $\set{C_{p,q}}_{p,q\in\bZ}$ of objects from $\A$ along with maps $d^h\colon C_{p,q}\to C_{p-1,q}$ and $d^v\colon C_{p,q}\to C_{p,q-1}$ (horizontal and
    vertical differentials) such that $d^h\circ d^h=d^v\circ d^v=d^v\circ d^h+d^h\circ d^v=0$.

    A bicomplex is {\emphcolor bounded} if $C$ has only finitely many nonzero components along every diagonal $p+q=n$ (e.g.\ if $C_{p,q}\neq0$ only when $p,q\geq0$).

\edefn

If $C_{pq}$ is a bicomplex, $C_{\bul q}$ and $C_{p\bul}$ are chain complexes for all $p,q$.

Note that due to anticommutativity ($d^vd^h=-d^hd^v$), $d^v$ are not maps $C_{\bul q}\to C_{\bul q-1}$, but they can be used to construct them: define
$$ f_{pq} = (-1)^pd^v_{pq}\colon C_{pq}\to C_{p,q-1} $$
Then $f_{\bul q}$ defines a chain map $C_{\bul q}\to C_{\bul q-1}$: $d^hf_{pq}=d^h(-1)^pd^v_{pq}=(-1)^{p-1}d^v_{p-1,q}d^h=f_{p-1,q}d^h$.
This ``sign trick'' allows us to identify the category of bicomplexes with $\Ch(\Ch(\A))$.

Let $C$ be a bicomplex, and define the chain complexes $\tot^\Pi(C)$ and $\tot^\oplus(C)$ by
$$ \tot^\Pi(C)_n = \prod_{p+q=n}C_{pq},\qquad \tot^\oplus(C)_n = \bigoplus_{p+q=n}C_{pq} $$
Then $d=d^h+d^v$ defines maps $\tot^\Pi(C)_n\to\tot^\Pi(C)_{n-1}$ and $\tot^\oplus(C)_n\to\tot^\oplus(C)_{n-1}$.
Explicitly, we map $x\in C_{pq}$ to $(d^hx,d^vx)\in C_{p-1,q}\times C_{p,q-1}$, or more suggestively:
$$ d\sum_{p+q=n}x_{pq} = \sum_{p+q=n}(d^hx_{pq}+d^vx_{pq}) $$
These are differential maps: $(d^h+d^v)(d^h+d^v)=0$ precisely because $d^hd^v+d^vd^h=0$.

And so $\tot^\Pi(C)$ and $\tot^\oplus(C)$ define chain complexes, called {\emphcolor total complexes}.
If $C$ is bounded, then $\tot^\Pi(C)=\tot^\oplus(C)$ (as there are only finitely many nonzero $C_{pq}$ for which $p+q=n$).

\bnote

    The total complexes do not exist in all Abelian categories!
    In order form them to exist, the infinite products and coproducts of $C_{pq}$ must exist.
    If $C$ is bounded, then they do, as the total complexes are equal to the finite biproducts of $C_{pq}$.
    But in general, we need $\A$ to be complete (equivalently have arbitrary products) to define $\tot^\Pi(C)$, and cocomplete (arbitrary coproducts) to define $\tot^\oplus(C)$.

    In particular the total complexes exist in $\Mod_R$ and $\Ch(\Mod_R)$.

\enote

\bdefn

    Let $C$ be a chain complex, and $n$ an integer.
    We define $\tau_{\geq n}C$ to be the subcomplex of $C$ defined by
    $$ (\tau_{\geq n}C)_i = \cases{0&if $i<n$\cr Z_n&if $i=n$\cr C_n&if $i>n$} $$
    This is called the {\emphcolor truncation} of $C$ below $n$.
    The quotient complex $\tau_{<n}C=C/\tau_{\geq n}C$ is the truncation of $C$ above $n$.

\edefn

Note that the $i$th boundary module and $i$th cycle module for $\tau_{\geq n}C$ are the same as those for $C$, except when $i<n$ when they are zero.
Thus
$$ H_i(\tau_{\geq n}C) = \cases{0&$i<n$\cr H_i(C)&$i\geq n$} $$
And so we get
$$ H_i(\tau_{<n}C) = \cases{H_i(C)&$i<n$\cr0&$i\geq n$} $$

The more naive method defining truncation does not give as clean of a homology group (defining the $n$th term to be $0$ or $C_n$, whatever).

\bdefn

    If $C$ is a chain or cochain complex and $p$ an integer we define a new complex (chain or cochain) $C[p]$ by
    $$ C[p]_n = C_{n+p}\qquad (\hbox{resp. }C[p]^n=C^{n-p}) $$
    with differential $(-1)^pd$.
    These are called the {\emphcolor $p$th translate} of $C$.
    The sign convention (multiplying by $(-1)^p$ will simplify notation later on).

\edefn

Clearly we have
$$ H_n(C[p]) = H_{n+p}(C)\qquad (\hbox{resp. }H^n(C[p])=H^{n-p}(C)) $$
We extand translation into a functor by shifting indices on a chain map: if $f\colon C\to D$ is a chain map then we define $f[p]$ by
$$ f[p]_n = f_{n+p}\qquad (\hbox{resp. }f[p]^n=f^{n-p}) $$

\bexam

    Let $C$ be a complex, then there is an exact sequence of complexes
    $$ 0 \to Z(C) \to C \xto{\ d\ } B(C)[-1] \to 0 $$
    Indeed, inclusion is mono and thus has trivial kernel; by definition the kernel of $d$ is the image of $Z(C)$.
    And the image of $d$ is by definition all boundaries.
    We must reindex for $d$ to be a chain map.

    Similarly we have an exact sequence
    $$ 0 \to H(C) \to C/B(C) \xto{\ d\ }Z(C)[-1] \to H(C)[-1] \to 0 $$

\eexam

\subsection{Long exact sequences}

Chain complexes are useful as they provide long exact sequences.
Our goal in this section is to prove the following:

\bthrm

    Let $0\to A\xto{\ f\ }B\xto{\ g\ }C\to0$ be a short exact sequence of chain complexes.
    Then there are natural maps $\partial_n\colon H_n(C)\to H_{n-1}(A)$ called {\emphcolor connecting homomorphisms} such that
    $$ \cdots\xto{\ g\ }H_{n+1}(C)\xto{\ \partial\ }H_n(A)\xto{\ f\ }H_n(B)\xto{\ g\ }H_n(C)\xto{\ \partial\ }H_{n-1}(A)\xto{\ f\ }\cdots $$
    is an exact sequence.

    Similarly if $0\to A\xto{\ f\ }B\xto{\ g\ }C\to0$ is a short exact sequence of cochain complexes, then there are natural maps $\partial_n\colon H^n(C)\to H^{n+1}(A)$
    and a long exact sequence:
    $$ \cdots\xto{\ g\ }H^{n-1}(C)\xto{\ \partial\ }H^n(A)\xto{\ f\ }H^n(B)\xto{\ g\ }H^n(C)\xto{\ \partial\ }H^{n+1}(A)\xto{\ f\ }\cdots $$

\ethrm

\bexam

    Note that if $0\to A\to B\to C\to 0$ is a short exact sequence such that two of the three complexes $A,B,C$ are exact (acyclic) then so is the third.
    If $A,C$ are exact then since we have an exact sequence $0=H_n(A)\to H_n(B)\to H_n(C)=0$, we must have $H_n(B)=0$ so $B$ is exact.
    The other cases are due to the above theorem: we will always have an exact sequence of the form $0\to H_n(X)\to0$ for $X=A,B,C$.

\eexam

We will require the following lemma:

\blemm[title=Snake lemma]

    Consider a commutative diagram of $R$-modules of the form:

    \centerline{\drawdiagram{
        &$A'$&$B'$&$C'$&$0$\cr
        $0$&$A$&$B$&$C$\cr
    }{
        \diagarrow{from={1,2},to={1,3},text=$\alpha$,y distance=-.25cm}
        \diagarrow{from={1,3},to={1,4},text=$p$,y distance=-.25cm}
        \diagarrow{from={1,4},to={1,5}}
        \diagarrow{from={1,2},to={2,2},text=$f$,x distance=-.25cm}
        \diagarrow{from={1,3},to={2,3},text=$g$,x distance=-.25cm}
        \diagarrow{from={1,4},to={2,4},text=$h$,x distance=-.25cm}
        \diagarrow{from={2,1},to={2,2}}
        \diagarrow{from={2,2},to={2,3},text=$i$,y distance=.25cm}
        \diagarrow{from={2,3},to={2,4},text=$\beta$,y distance=.25cm}
    }}

    If the rows are exact, then there is an exact sequence
    $$ 0\to\ker f\to\ker g\to\ker h\xto{\ \partial\ }\coker f\to\coker g\to\coker h\to0 $$
    where $\partial$ is defined by $\partial(c)=i^{-1}gp^{-1}(c)$ for $c\in\ker h$.
    Moreover, if $\alpha\colon A'\to B'$ is monic then so is $\ker f\to\ker g$; if $\beta\colon B\to C$ is onto then so is $\coker g\to\coker h$.

\elemm

\bproof

    The top row is exact, so $p$ is surjective.
    Thus for $c\in\ker h$ there is an $x\in B'$ such that $p(x)=c$.
    Note then that $\beta gx=hpx=hc=0$ and so $gx\in\ker\beta=\im i$ so $i^{-1}gx$ exists.
    But this doesn't ensure that this is well-defined: if $p(x)=p(y)=c$ we still need to show that $i^{-1}gx=i^{-1}gy$.

    Note that $i^{-1}$ here maps into the cokernel of $f$, that is to say we map to $i^{-1}\bul+\im f$.
    So a more exact way of writing $\partial$ would be by
    $$ \partial(c) = i^{-1}gp^{-1}(c) + \im f $$

    Suppose $px=0$ then $x\in\ker p=\im\alpha$ so $x=\alpha(a)$ for $a\in A'$.
    Then we have $gx=ifa$ and so $i^{-1}gx=fa\in\im f$, and so $i^{-1}gx=0$ in $\coker f$.
    Thus if $p(x)=p(y)$ then $p(x-y)=0$ and so $i^{-1}gx-i^{-1}gy\in\im f$ meaning they define the same coset in $\coker f$; $\partial$ is well-defined.
    It is also clearly a homomorphism, as if $p(x)=c$ and $p(y)=c'$ then $\partial(c+c')=i^{-1}gx+i^{-1}gy=\partial(c)+\partial(c')$.

    Now $0\to\ker f\xto\alpha\ker g\xto p\ker h$ is an exact sequence, and so is $\coker f\xto i\coker g\xto\beta\coker h\to0$.
    Adding $\partial$ allows us to connect these two exact sequences.
    \qed

\eproof

We can extend the snake lemma to arbitrary Abelian categories.
The Freyd-Mitchell embedding theorem allows us to embed small Abelian categories into categories of modules.
So take the Abelian subcategory generated by the commutative, embed it into a category of modules, and the result is given (details are omitted).

Now we construct the natural maps $\partial\colon H_n(C)\to H_{n-1}(A)$ in our goal theorem, associated with $0\to A\to B\to C\to0$.
Consider the commutative diagram (we write $dA_{n+1}$ for $B_n(A)$ since we already have a chain complex $B_n$)

\centerline{\drawdiagram{
    &$A_n/dA_{n+1}$&$B_n/dB_{n+1}$&$C_n/dC_{n+1}$&$0$\cr
    $0$&$Z_{n-1}(A)$&$Z_{n-1}(B)$&$Z_{n-1}(C)$
}{
    \diagarrow{from={1,2},to={1,3}}
    \diagarrow{from={1,3},to={1,4}}
    \diagarrow{from={1,4},to={1,5}}
    \diagarrow{from={2,1},to={2,2}}
    \diagarrow{from={2,2},to={2,3}}
    \diagarrow{from={2,3},to={2,4}}
    \diagarrow{from={1,2},to={2,2},text=$d_A$,x distance=-.25cm}
    \diagarrow{from={1,3},to={2,3},text=$d_B$,x distance=-.25cm}
    \diagarrow{from={1,4},to={2,4},text=$d_C$,x distance=-.25cm}
}}

The rows are exact.
The kernel of $d_A$ is $H_n(A)$, its cokernel is $H_{n-1}(A)$.
Thus the snake lemma yields an exact sequence
$$ H_n(A) \to H_n(B) \to H_n(C) \xto{\ \partial\ } H_{n-1}(A) \to H_{n-1}(B) \to H_{n-1}(C) $$
pasting these sequences together gives us the desired long exact sequence.
\qed

Using the snake lemma's construction, we can understand $\partial$.
Consider the short exact sequence
$$ 0 \to A \xto{\ f\ } B \xto{\ g \ } C \to 0 $$
then $\partial(z)=f^{-1}dg^{-1}(z)$.
That is, for $[z]\in H_n(C)$ consider its lift $b\in B_n$ (meaning $g(b)=z$), apply $d$ and then $f^{-1}$: $\partial(z)=f^{-1}(db)$.

Now $\partial$ is natural in the following sense: consider the category $\c S$ whose objects are short exact sequences of chain complexes in $\A$, and whose morphisms commute with these.
Similarly there is a category $\c L$ of long exact sequences in $\A$.

\bprop

    The long exact sequence is a functor from $\c S$ to $\c L$.

\eprop

We omit the proof here; it is a direct result of the construction of the connecting morphisms (from the snake lemma) and definitions.

\bexam

    Suppose $f\colon C\to D$ is a morphism of chain complexes such that $\ker f,\coker f$ are acylic.
    Then $f$ is a quasi-isomorphism.
    Indeed, we have an exact sequence
    $$ \ker f \to C \xto{\ f\ } D \to \coker f $$
    This gives us an exact sequence of homology modules
    $$ 0 = H_n(\ker f) \to H_n(C) \xto{\ f\ } H_n(D) \to H_n(\coker f) = 0 $$
    meaning that the induced morphism between homology modules must be an isomorphism: $f$ is quasi-isomorphic.

\eexam

\bexam

    Let $0\to A\xto f B\xto g C\to0$ be a short exact sequence of bicomplexes.
    This gives us a short exact sequence of total complexes.
    Indeed, $f$ induces a map $\tot(A)\to\tot(B)$ by mapping $(a_{pq})$ to $(f_{pq}a_{pq})$; the induced maps still form into a short exact sequence.
    Thus if $\tot(A)$ or $\tot(C)$ are acylic, then the other two are quasi-isomorphic (since the exact sequence of homology modules is of the form $0\to X\to Y\to0$).

\eexam

\subsection{Chain homotopies}

Recall that chain complexes of vector spaces over a field decompose as $C_n=Z_n\oplus B'_n$ where $B'_n\cong C_n/Z_n\cong d(C_n)=B_{n-1}$, and $Z_n=B_n\oplus H'_n$ where $H'_n\cong Z_n/B_n=H_n(C)$.
So we have compositions
$$ C_n \xto{\rm proj} Z_n \xto{\rm proj} B_n \cong B'_{n+1} \subseteq C_{n+1} $$
giving us splitting maps $s_n\colon C_n\to C_{n+1}$ such that $d=dsd$:
$$ C_{n+1} \xto{\ d\ } C_n \to Z_n \to B_n \cong B'_{n+1} \subseteq C_{n+1} \xto{\ d\ } C_n $$
the map $C_{n+1}\to B_n$ simply maps $a$ to $da$.
The isomorphism $B_n\cong B'_{n+1}$ maps $da$ to $a$'s projection onto $B'_{n+1}$ (orthogonal to $Z'_{n+1}$; following the isomorphisms), which $d$ then maps to $da$ (since $da$ is equal to $d$ on the projection).
Thus $dsd=d$, as desired.

The composition $ds$ is the projection of $C_n$ onto $B_n$, and $sd$ the projection onto $B'_n$.
Thus $ds+sd$ is an endomorphism of $C_n$ whose kernel is $H'_n$, isomorphic to the homology $H_n(C)$.
Its image is $B_n\oplus B'_n$ and as such its cokernel is isomorphic to $H_n(C)$ as well.
So the kernel and cokernel of $ds+sd$ are isomorphic to the homology complex $H_\bul(C)$ (the differentials here are zero, so the homology modules are the same as $C$'s).
And so the chain maps $H_\bul(C)\to C$ and $C\to H_\bul(C)$ are quasi-isomorphisms.
Furthermore, since the (co)kernel of $ds+sd$ is the homology complex, $C$ is exact (acylic) iff $ds+sd$ is an isomorphism.

This is not true over an arbitrary ring $R$, so we give this a name.

\bdefn

    A complex $C$ is called {\emphcolor split} if there are maps $s_n\colon C_n\to C_{n+1}$ such that $d=dsd$.
    These maps $s_n$ are called {\emphcolor splitting maps}.
    If $C$ is further acylic, then $C$ is said to be {\emphcolor split exact}.

\edefn

\bprop

    A chain complex $C$ is split iff $C_n\cong Z_n\oplus B'_n$ with $Z_n=B_n\oplus H'_n$.

\eprop

\bproof

    If $C$ is split then define $B'_n$ to be the image of $s_{n-1}d_n$.
    Since $dsd=d$, we have that $\ker sd=\ker d=Z_n$ and $sdsd=sd$, so $sd$ is a projection operator whose image is $B'_n$ and whose kernel is $Z_n$, thus $C_n=Z_n\oplus B'_n$ as desired.
    Similarly, $ds$ is a projection operator whose image is $B_n$; restricting it to $Z_n$ gives us $Z_n=B_n\oplus H'_n$.

    If $C_n\cong Z_n\oplus B'_n$ and $Z_n=B_n\oplus H'_n$ then we have $B'_n\cong B_{n-1}$.
    And so we can map $C_n\to C_{n+1}$ along $B_n\cong B'_{n+1}$; $(x,y,z)\in B_n\oplus H'_n\oplus B'_n$ maps to $(0,0,x)\in B_{n+1}\oplus H'_{n+1}\oplus B'_{n+1}$.
    Then $dsd$ maps $(x,y,z)\in C_{n+1}$ to $(z,0,0)\in C_n$ to $(0,0,z)\in C_{n+1}$ to $(z,0,0)\in C_n$, which is the same as $d(x,y,z)$; $dsd=d$.
    \qed

\eproof

\bprop

    A chain complex $C$ is split exact iff there exist maps such that $ds+sd$ is the identity.

\eprop

\bproof

    If $ds+sd={\rm id}$ then we get that $(ds+sd)d=d$, but the left-hand side is $dsd+sd^2=dsd$, so $C$ is split.
    And $C$ is split then we showed that $C_n\cong B_n\oplus H'_n\oplus B'_n$ with $ds,sd$ the projections onto $B_n,B'_n$ respectively.
    $C$ is then split exact iff $H_n(C)\cong H'_n=0$, which is iff the projection onto $B_n\oplus B'_n$ is the identity.
    This is $ds+sd$, and so $C$ is split exact iff $ds+sd$ is the identity, as desired.
    \qed

\eproof

\bexam

    Consider the complex
    $$ \cdots\xto{\ \cdot2\ }\bZ/4\bZ\xto{\ \cdot2\ }\bZ/4\bZ\xto{\ \cdot2\ }\cdots $$
    This is an acylic complex: the image and kernel of multiplication by $2$ are $\set{0,2}$.
    But it is not split exact: $ds+sd$ can never be the identity.

\eexam

Now suppose we have two chain complexes $C,D$ and maps $s_n\colon C_n\to D_{n+1}$.
Define $f_n\colon C_n\to D_n$ by $f_n=d_{n+1}s_n+s_{n-1}d_n$: $f=ds+sd$.
Then notice that
$$ df = d(ds+sd) = dsd = (ds+sd)d = fd $$
so $f=ds+sd$ is a chain map $C\to D$.

\bdefn

    A chain map $f\colon C\to D$ is {\emphcolor null homotopic} if there are maps $s_n\colon C_n\to D_{n+1}$ such that $f=ds+sd$.
    The maps $s_n$ are called a {\emphcolor chain contraction} of $f$.

\edefn

Immediately we get:

\bprop

    $C$ is split exact iff the identity is null homotopic.

\eprop

Since the sum of chain maps is a chain map, the sum of any chain map with a null homotopic map is again a chain map.

\bdefn

    We say that two maps $f,g\colon C\to D$ are {\emphcolor chain homotopic} if their difference is null homotopic.
    If $f-g=sd+ds$ then the maps $s_n$ are called a {\emphcolor chain homotopy} from $f$ to $g$.
    Two maps $f\colon C\tof D\colon g$ form a {\emphcolor chain homotopy equivalence} if their compositions are chain homotopic to the respective identities.
    $C$ and $D$ are {\emphcolor chain homotopy equivalent} if there exists a chain homotopy equivalence between them.

\edefn

In particular $f$ is null homotopic iff it is chain homotopic to the zero map.

\blemm

    If $f\colon C\to D$ is null homotopic then the induced maps $f\colon H_n(C)\to H_n(D)$ are zero.
    And so if $f,g$ are chain homotopic their induced maps are equal.

\elemm

\bproof

    The latter claim follows from the first: the induced map $H_n(f-g)$ is equal to $H_n(f)-H_n(g)$.
    So if $f-g$ is null homotopic, then $H_n(f)-H_n(g)=0$.

    So suppose $f=ds+sd$, then for $a\in Z_n(C)$ we have $fa=dsa+sda=dsa\in B_n(D)$.
    Thus $f$ maps cycles to boundaries, and so its induced map is zero.
    \qed

\eproof

\bexam

    If $C$ is split, then the chain complex $H_\bul(C)$ and $C$ are chain homotopy equivalent.
    Indeed, consider $f\colon C\to C$ given by $f=\id-ds-sd$ (the difference between a chain map and a null homotopy).
    Then $df=0$, so $f$ maps into $Z_n$, that is it defines a map $f\colon C\to H_\bul(C)$ (alternatively recall that $C_n\cong B_n\oplus H_n\oplus B'_n$, this is the projection onto $H_n$).
    Then define $g\colon H_\bul(C)\to C$ by inclusion (recalling that $C_n\cong B_n\oplus H_n\oplus B'_n$).

    Then $fg$ is just the identity on $H_\bul(C)$.
    And $gf$ is projection onto $H_n$; it is $\id-ds-sd$ which is homotopy equivalent to $\id$.
    Thus $f,g$ form a chain homotopy equivalence.
    \qed

\eexam

Notice that chain homotopy equivalence is an equivalence relation: it is clearly reflexive (taking $s=0$) and symmetric (taking $-s$).
And if $f-g=sd+ds$ and $g-h=s'd+ds'$ then $f-h=(s+s')d+d(s+s')$, so it is transitive.
This equivalence relation preserves addition: if $f-f'=sd+ds$ and $g-g'=s'd+ds'$ then $(f+g)-(f'+g')=(s+s')d+d(s+s')$.

Let us construct the category $\K$ whose objects are chain complexes and whose morphisms $C\to D$ are equivalence classes of chain maps.
We need to show that composition is well-defined: if $f-g=sd+ds$ and $u$ is another chain map then $uf-ug=usd+uds=usd+dus$, so $uf$ and $ug$ are homotopic (similar for $fu,gu$).
Thus if $f\sim g$ and $u\sim v$ we have $fu\sim gu\sim gv$, so composition is well-defined.

This turns $\K$ into an additive category; it has biproducts.
It is not an Abelian category as it doesn't necessarily have kernels.

Since we showed that chain homotopic maps have the same image under homology functors, homology gives functors $\K\to\A$.
In particular chain homotopic equivalent complexes have isomorphic homology (since homotopy equivalence is isomorphism in $\K$):

\bthrm

    If $C$ and $D$ are chain homotopic equivalent, then $H_\bul(C)\cong H_\bul(D)$.

\ethrm

\subsection{Mapping cones and cylinders}

\bdefn

    Let $f\colon B\to C$ be a map of chain complexes.
    The {\emphcolor mapping cone} of $f$ is the chain complex $\cone(f)$ whose $n$th component is $B_{n-1}\oplus C_n$.
    The differential is given by $d(b,c)=(-d(b),d(c)-f(b))$.

    For cochain complexes, the $n$th component is $B^{n+1}\oplus C^n$, and the differential has the same formula.

\edefn

We can represent the differential by the matrix $\pmatrix{-d_B&0\cr-f&+d_C}$.

\bexam

    Let $\cone(C)$ denote the mapping cone of the identity map $\id_C$; its differential is given by $d(b,c)=(-db,dc-b)$.
    Then $s(b,c)=(-c,0)$ is a splitting map $\cone(C)_n=C_{n-1}\oplus C_n\to C_n\oplus C_{n+1}=\cone(C)_{n+1}$:
    $$ (ds+sd)(b,c) = d(-c,0) + s(-db,dc-b) = (dc,c) + (b-dc,0) = (b,c) $$
    so $\cone(C)$ is split exact.

\eexam

\bprop

    Let $f\colon C\to D$ be a map of complexes.
    Then $f$ is null homotopic iff it extends to a map $(-s,f)\colon\cone(C)\to D$.

\eprop

\bproof

    Note that $(-s,f)$ is a chain map iff $d(-s,f)(b,c)=(-s,f)d(b,c)$ for all $b,c$.
    These are equal to
    $$ d(-s,f)(b,c) = -ds(b) + df(c),\qquad (-s,f)d(b,c) = (-s,f)(-db,dc-b) = sd(b) + fd(c) - f(b) $$
    since $f$ is a chain map $fd=df$, and so these are equal iff
    $$ -ds(b) = sd(b) - f(b) $$
    i.e.\ iff $f=sd+ds$, as desired.
    \qed

\eproof

Consider a chain map $f\colon H_\bul(B)\to H_\bul(C)$, then there is a short exact sequence
$$ 0 \to C \xto{\ \iota\ } \cone(f) \xto{\ \delta\ } B[-1] \to 0 $$
where the left map maps $c\in C_n$ to $(0,c)\in\cone(f)_n$, and the right maps $(b,c)$ to $-b$.
These are chain maps: $d(0,c)=(0,dc)$ and $d(b,c)=(-db,dc-fb)\mapsto -db$.
The left map is an injection, the right a surjection, and the image of the left is the kernel of the second (the component $C_n$).

And so this gives us a long exact sequence (recalling that $H_{n+1}(B[-1])=H_n(B)$)
$$ \cdots\to H_{n+1}(\cone f) \xto{\ \delta\ } H_n(B) \xto{\ \partial\ } H_n(C) \to H_n(\cone f) \xto{\ \delta\ } H_{n-1}(B) \to \cdots $$
Let us understand the connecting maps $\partial$.
Recall that $\partial=\iota^{-1}d_{\cone(f)}\delta^{-1}$, so for a cycle $b\in B_n$ we can lift it to $(-b,0)\in\cone(f)$.
Then applying the differential gives $(db,fb)=(0,fb)$ which maps to $fb$.
Thus we have $\partial([b])=f([b])$ for all $[b]\in H_\bul(B)$, so our connecting map is simply $f$.

We have proven the following.

\bprop

    For any map $f\colon H_\bul(B)\to H_\bul(C)$ we have a long exact sequence
    $$ \cdots\to H_{n+1}(\cone f) \to H_n(B) \xto{\ f\ } H_n(C) \to H_n(\cone f) \to H_{n-1}(B) \xto{\ f\ } \cdots $$

\eprop

\bcoro

    $f\colon B\to C$ is a quasi-isomorphism iff $\cone f$ is acyclic.

\ecoro

\bproof

    If $\cone f$ is acyclic the above long exact sequence gives us exact sequences $0\to H_n(B)\xto f H_n(C)\to0$, so the induced map is an isomorphism.
    And if $f$ is a quasi-isomorphism, we must have that $\pi$ and $\delta$ are zero on the homology modules, since $\ker\delta=\im\pi=0$ we must have that $H_n(\cone f)=0$, so $\cone f$ is acyclic.
    \qed

\eproof

We have a similar construction:

\bdefn

    Let $f\colon B\to C$ be a chain map, then define the {\emphcolor mapping cylinder} of $f$ to be the complex $\cyl(f)$ whose $n$th component is $B_n\oplus B_{n-1}\oplus C_n$.
    The differential is
    $$ d(b,b',c) = (d(b)+b',-d(b'),d(c)-f(b')) $$

\edefn

We can represent the differential as the matrix
$$ d = \pmatrix{d_B&\id_B&0\cr0&-d_B&0\cr0&-f&d_C} $$
We see that this is a complex:
$$ d^2 = \pmatrix{d_B^2&d_B-d_B&0\cr0&d_B^2&0\cr0&fd_B-d_Cf&d_C^2} = 0 $$

\bprop

    Let $\cyl(C)$ be the mapping cylinder of $\id_C$.
    Then $f,g\colon C\to D$ are chain homotopic iff they extend to a map $(f,s,g)\colon\cyl(C)\to D$.

\eprop

\bproof

    $(f,s,g)$ is a chain map iff $d(f,s,g)(b,b',c)=(f,s,g)d(b,b',c)$ for all $b,b',c$.
    These are equal to
    $$ d(f,s,g)(b,b',c) = df(b) + ds(b') + dg(c),\qquad (f,s,g)d(b,b',c) = (f,s,g)(d(b)+b',-d(b'),d(c)-b') = fd(b) + f(b') - sd(b') + gd(c) - g(b') $$
    These are equal iff $f(b')-sd(b')-g(b')=ds(b')$ i.e.\ iff $f-g=ds+sd$, as desired.
    \qed

\eproof

\blemm

    The subcomplex whose elements are of the form $(0,0,c)$ is isomorphic to $C$ and the corresponding inclusion $\alpha\colon C\to\cyl(f)$ is a quasi-isomorphism.

\elemm

\bproof

    Clearly $\alpha\colon c\mapsto(0,0,c)$ is an injection whose image is this subcomplex.
    We need to show that $\alpha$ is a chain map:
    $$ d\alpha(c) = d(0,0,c) = (0,0,dc) = \alpha(dc) $$
    as desired.

    Now, up to isomorphism, $\cyl(f)/\alpha(C)$ has components of the form $B_n\oplus B_{n-1}$ and the differential is given by $d(b,b')=(d(b)+b',-d(b'))$.
    This is clearly isomorphic to $\cone(-\id_B)$, which we showed is split exact (we showed $\cone(\id_B)$ is).
    So we have a short exact sequence
    $$ 0 \to C \xto{\ \alpha\ } \cyl(f) \to \cone(-\id_B) \to 0 $$
    and the homology sequence gives us $0=H_{n+1}(\cone(-\id_B))\to H_n(C)\xto\alpha H_n(\cyl(f))\to H_n(\cone(-\id_B))=0$, that is to say $\alpha$ is a quasi-isomorphism.
    \qed

\eproof

We can actually strengthen this:

\blemm

    $\alpha$ is actually a chain homotopy equivalence (thus an isomorphism in $\K$, so in particular a quasi-isomorphism).

\elemm

\bproof

    Define $\beta\colon\cyl(f)\to C$ by $\beta(b,b',c)=f(b)+c$.
    This is a chain map:
    $$ d\beta(b,b',c) = df(b) + dc,\qquad \beta d(b,b',c) = \beta(d(b)+b',-d(b'),d(c)-f(b')) = f(db+b') + dc - f(b') = fd(b) + dc $$
    Now $\beta\alpha=\id_C$ and defining $s_n\colon\cyl(f)_n\to\cyl(f)_{n+1}$ by $s(b,b',c)=(0,b,0)$ gives a chain homotopy from $\alpha\beta$ to $\id$:
    $$ ds(b,b',c) + sd(b,b',c) = d(0,b,0) + s(d(b)+b',-d(b'),d(c)-f(b')) = (b,-d(b),-f(b)) + (0,d(b)+b',0) = (b,b',-f(b)) $$
    and
    $$ \alpha\beta(b,b',c) = (0,0,f(b)+c) $$
    thus $\id-\alpha\beta=ds+sd$ as desired.
    \qed

\eproof

To summarize, $\alpha\colon C\tof\cyl(C)\colon\beta$ form a chain homotopy equivalence where
$$ \alpha(c) = (0,0,c),\qquad \beta(b,b',c) = f(b) + c $$

Now consider the map $\iota\colon B\to\cyl(f)$ by $\iota(b)=(b,0,0)$.
This is a chain map:
$$ d\iota(b) = d(b,0,0) = (db,0,0) = \iota(db) $$
and so the subcomplex of $\cyl(f)$ whose elements are of the form $(b,0,0)$ is isomorphic to $B$.
And the quotient complex $\cyl(f)/\iota(B)$ is isomorphic to $\cone(f)$: its terms are of the form $B_{n-1}\oplus C_n$ and the differential is given by $d(b',c)=(-d(b'),d(c)-f(b'))$.

Consider then the composite $B\xto{\ \iota\ }\cyl(f)\xto{\ \beta\ }C$: this is just $f$.
Thus homology $f\colon H(B)\to H(C)$ factors through $H(B)\to H(\cyl(f))$.
Then we can construct a commutative diagram:

\centerline{\drawdiagram{
    &&$C$\cr
    $0$&$B$&$\cyl(f)$&$\cone(f)$&$0$\cr
    &$0$&$C$&$\cone(f)$&$B[-1]$&$0$
}{
    \diagarrow{from={2,1}, to={2,2}}
    \diagarrow{from={2,2}, to={2,3}}
    \diagarrow{from={2,3}, to={2,4}}
    \diagarrow{from={2,4}, to={2,5}}
    %
    \diagarrow{from={3,2}, to={3,3}}
    \diagarrow{from={3,3}, to={3,4}}
    \diagarrow{from={3,4}, to={3,5}, text=$\delta$, y distance=.25cm}
    \diagarrow{from={3,5}, to={3,6}}
    %
    \diagarrow{from={2,2}, to={1,3}, text=$f$, x distance=-.25cm}
    \diagarrow{from={2,3}, to={1,3}, text=$\beta$, x distance=.25cm}
    \diagarrow{from={3,3}, to={2,3}, text=$\alpha$, x distance=.25cm}
    \diagarrow{from={3,4}, to={2,4}, text={$=$}, x distance=.25cm, left cap=<}
}}

The rows are exact (the map $\cyl(f)\to\cone(f)$ is given by the isomorphism $\cyl(f)/\iota(B)\cong\cone(f)$).
Thus we have long exact sequences (recall that the connecting morphisms of the bottom homology are the induced maps of $f$):

\bigskip
\centerline{\drawdiagram{
    $\cdots$&$H_n(B)$&$H_n(\cyl(f))$&$H_n(\cone(f))$&$H_{n-1}(B)$&$\cdots$\cr
    $\cdots$&$H_{n+1}(B[-1])$&$H_n(C)$&$H_n(\cone(f))$&$H_n(B[-1])$&$\cdots$
}{
    \diagarrow{from={1,1}, to={1,2}, text=$-\partial$, y distance=.25cm}
    \diagarrow{from={1,4}, to={1,5}, text=$-\partial$, y distance=.25cm}
    \diagarrow{from={1,2}, to={1,3}}
    \diagarrow{from={1,3}, to={1,4}}
    \diagarrow{from={1,5}, to={1,6}}
    %
    \diagarrow{from={2,2}, to={2,3}, text=$f$, y distance=-.25cm}
    \diagarrow{from={2,5}, to={2,6}, text=$f$, y distance=-.25cm}
    \diagarrow{from={2,4}, to={2,5}, text=$\delta$, y distance=-.25cm}
    \diagarrow{from={2,1}, to={2,2}}
    \diagarrow{from={2,3}, to={2,4}}
    %
    \diagarrow{from={1,2}, to={2,2}, text={$=$}, x distance=.25cm, left cap=<}
    \diagarrow{from={1,3}, to={2,3}, text=$\cong$, x distance=.25cm, left cap=<}
    \diagarrow{from={1,4}, to={2,4}, text={$=$}, x distance=.25cm, left cap=<}
    \diagarrow{from={1,5}, to={2,5}, text={$=$}, x distance=.25cm, left cap=<}
    \diagarrow{from={1,2}, to={2,3}, text=$f$, x distance=.25cm}
}}

We take the inverse of the top connecting map precisely because we want this diagram to commute.

To check that this diagram commutes we need only check that the right square with $-\partial$ and $\delta$ does.
$\partial$ works as follows: let $[(b,c)]$ be a cycle in $\cone(f)$, so $db=0$ and $dc=fb$.
Then we lift it to $\cyl(f)$ to $(0,b,c)$ and then apply the differential $d(0,b,c)=(b,-db,dc-fb)=(b,0,0)$.
Thus $\partial$ maps $[(b,c)]\in H(B)$ to $[b]\in H_{n-1}(B)$.
$\delta$ maps $[(b,c)]$ to $[-b]$, i.e.\ $\delta=-\partial$ as desired.

And so these diagrams are isomorphic.

Now suppose we have a short exact sequence $0\to B\xto fC\xto gD\to0$.
We define a chain map $\phi\colon\cone(f)\to D$ by $\phi(b,c)=g(c)$:
$$ \phi d(b,c) = \phi(-db,dc-fb) = gdc - gfb = gdc = d\phi(b,c) $$
Then we have long exact sequences

\bigskip
\centerline{\drawdiagram{
    $0$&$B$&$C$&$D$&$0$\cr
    $0$&$B$&$\cyl(f)$&$\cone(f)$&$0$\cr
    &$0$&$C$&$\cone(f)$&$B[-1]$&$0$
}{
    \diagarrow{from={1,1}, to={1,2}}
    \diagarrow{from={1,2}, to={1,3}, text=$f$, y distance=.25cm}
    \diagarrow{from={1,3}, to={1,4}, text=$g$, y distance=.25cm}
    \diagarrow{from={1,4}, to={1,5}}
    %
    \diagarrow{from={2,1}, to={2,2}}
    \diagarrow{from={2,2}, to={2,3}}
    \diagarrow{from={2,3}, to={2,4}}
    \diagarrow{from={2,4}, to={2,5}}
    %
    \diagarrow{from={3,2}, to={3,3}}
    \diagarrow{from={3,3}, to={3,4}}
    \diagarrow{from={3,4}, to={3,5}, text=$\delta$, y distance=.25cm}
    \diagarrow{from={3,5}, to={3,6}}
    %
    \diagarrow{from={1,2}, to={2,2}, text={$=$}, x distance=.25cm, left cap=<}
    \diagarrow{from={2,3}, to={1,3}, text=$\beta$, x distance=.25cm}
    \diagarrow{from={3,3}, to={2,3}, text=$\alpha$, x distance=.25cm}
    \diagarrow{from={3,4}, to={2,4}, text={$=$}, x distance=.25cm, left cap=<}
    \diagarrow{from={2,4}, to={1,4}, text=$\phi$, x distance=.25cm}
}}

Since our map from short to long exact sequences is a functor, and furthermore from the $5$-lemma (a special case saying that if the rows of a two-row five-column commutative diagram are exact,
and the outer four vertical maps are isomorphisms, then so is the innermost) $\phi$ is a quasi-isomorphism, since $\beta$ is.

Thus the rows all generate isomorphic long exact sequences (though we must invert the connecting morphism of the top row, as per the previous argument).

\bprop

    Let $f\colon B\to C$ be a chain map, and let $v\colon C\to\cone(f)$ be the inclusion $c\mapsto(0,c)$.
    There is a chain homotopy equivalence $\cone(v)\to B[-1]$.

\eprop

\bproof

    The $n$th component of $\cone(v)$ is $C_{n-1}\oplus B_{n-1}\oplus C_n$ and its differential is
    $$ d(c',b,c) = (-d(c'),d(b,c)-v(c')) = (-d(c'), -d(b), d(c) - f(b) - c') $$
    The maps $\alpha\colon\cone(v)\tof B[-1]\colon\beta$ given by $\alpha(c',b,c)=b$ and $\beta(b)=(f(b),b,0)$ form a chain homotopy equivalence.
    $\alpha\beta=\id$ and $\beta\alpha(c',b,c)=(-f(b),b,0)$ is chain homotopic to $\id$.
    Define $s_n\colon C_{n-1}\oplus B_{n-1}\oplus C_n\to C_n\oplus B_n\oplus C_{n+1}$ by $s_n(c',b,c)=(-c,0,0)$.
    Then
    $$ ds(c',b,c) + sd(c',b,c) = d(-c,0,0) + s(-dc',-db,dc-fb-c') = (dc,0,c) + (c'+fb-dc,0,0) = (c'+fb,0,c) $$
    which is equal to $(c',b,c)-\beta\alpha(c',b,c)=(c'+fb,0,c)$, meaning $\id$ and $\beta\alpha$ are chain homotopic.
    \qed

\eproof

\bexam

    Suppose $C$ and $C'$ are split complexes with splitting maps $s,s'$.
    Then $\sigma(c,c')=(-sc,s'c'-s'fsc)$ is a splitting map of $\cone(f)$ for $f\colon C\to C'$ iff the induced map $H(f)$ is zero.
    Indeed:
    $$ d\sigma d(c,c') = d\sigma(-dc,dc'-fc) = d(sdc, s'dc'-s'fc+s'fsdc) = (-dc, ds'dc' - ds'fc + ds'fsdc - fsdc) $$
    Since $s'$ is split, $ds'd=d$ so this is equal to $(-dc,dc'-(ds'f-ds'fsd+fsd)c)$, which is equal to $d$ iff $ds'f-ds'fsd+fsd=f$.
    Now since $C$ is split, we can embed $H(C)$ into it with projection $\id-ds-sd$, and so $H(f)$ is zero iff $f=fds+fsd$.
    So if $H(f)=0$ then we have
    $$ ds'f - ds'fsd + fsd = ds'(f-fsd) + fsd = ds'fds + fsd = dfs + fsd = f $$
    (since $fd=df$) as desired.
    Now if $\sigma$ is split then for $z\in Z_n(C)$ we have that $fz=ds'fz\in B_n(C')$, so $f$ is zero on the homology modules.

\eexam

\subsection{More on Abelian categories}

We begin by mentioning an important theorem in the study of Abelian categories, without proof:

\bthrm[title=Freyd-Mitchell]

    If $\A$ is a small Abelian category, then there exists an exact fully faithful functor from $\A$ to $\Mod_R$ for some ring $R$.

\ethrm

What is an {\it exact\/} functor?
We will define this soon.

\blemm

    Let $\C\subseteq\A$ be a full subcategory of the Abelian category $\A$.
    \benum
        \item $\C$ is additive iff $\C$ has $\A$'s zero objects and is closed under $\A$'s biproducts.
        \item $\C$ is itself Abelian and an exact subcategory of $\A$ iff $\C$ is additive and closed under $\ker,\coker$.
    \eenum

\elemm

\bexam

    \benum
        \item $\Mod_R$ is Abelian and the subcategory of finitely-generated $R$-modules forms an additive category.
        It is Abelian iff $R$ is Noetherian.
        \item $\Ab$ is Abelian, and the subcategory of torsion-free Abelian groups is an additive subcategory.
        The subcategory of $p$-groups is Abelian, so is the subcategory of finite $p$-groups.
    \eenum

\eexam

\bprop

    Let $\A$ be an Abelian category and $\C$ an arbitrary category.
    The {\emphcolor functor category} $[\C,\A]$ is Abelian.

\eprop

\bproof

    The zero object in $[\C,\A]$ is the constant functor to $0\in\A$.
    For two functors $F,G\colon\C\to\A$ define $F\oplus G\colon\C\to\A$ to map $F\oplus G(c)=Fc\oplus Gc$ and $f\colon c\to d$ to the matrix $\diag(Ff,Gf)$.
    This is a biproduct.
    The kernel of $\alpha\colon F\To G$ is $(\ker\alpha)(c)=\ker(\alpha(c))$, similar for cokernels.
    \qed

\eproof

\bdefn

    For a topological space $X$, consider the partial order $\U$ of open subsets of $X$, considered as a category.
    A {\emphcolor presheaf} (with values in $\A$) is a contravariant functor $F\colon\U^\op\to\A$ such that $F\emptyset=0$.
    The full subcategory of $[\U^\op,\A]$ of presheaves is denoted $\Psh_\A(X)$ (some authors do not require $F\emptyset=0$, in which case the category of presheaves is just the functor category).
    This is an Abelian category.

\edefn

An example of a presheaf with values in $\Mod_\bR$ is the functor $C^0$ mapping an open set $U$ to the set $C^0(U)$ of continuous functions $U\to\bR$.
For an inclusion $U\subseteq V$, the morphism $C^0(V)\to C^0(U)$ is given by restricting functions to $U$.

For a presheaf $F$, for $U\subseteq V$ we call the induced morphism $FV\to FU$ a {\emphcolor restriction}, and for $f\in FV$ its image under this morphism is also denoted $f|_U$.

\bdefn

    A {\emphcolor sheaf} on a topological space $X$ is a presheaf $F$ which satisfies the following {\emphcolor sheaf axiom}:
    let $\set{U_i}$ be an open cover of $U$, an open subset of $X$.
    If $\set{f_i\in F(U_i)}$ are such that $f_i$ and $f_j$ agree on $U_i\cap U_j$ (i.e.\ their restrictions onto $U_i\cap U_j$ are equal), then there is a unique $f\in F(U)$ whose restriction to
    $U_i$ is $f_i$.

\edefn

Note by uniqueness if $f|_{U_i}=g|_{U_i}$ for all $i$ then $f=g$.
The sheaf axiom is equivalent to saying that if $f\in F(U)$ vanishes on every $U_i$, then $f=0$.
Thus we can write the sheaf axiom in element-free form.
It is equivalent to saying that for any open $U$ and open cover $\set{U_i}_i$ of $U$, the following sequence is exact:
$$ 0 \to F(U) \to \prod_i F(U_i) \to \prod_{i,j}F(U_i\cap U_j) $$
the first map maps $f\in F(U)$ to $(f|_{U_i})_i$, the second maps $(f_i)_i$ to $(f_i|_{U_i\cap U_j}-f_j|_{U_i\cap U_j})_{i,j}$.

\bexam

    Let $A$ be an Abelian group.
    For any open subset $U$ of $X$, consider $A(U)$ the set of continuous maps from $U$ to the discrete topological space on $A$.
    $A(V)\to A(U)$ given by restriction forms this into a sheaf.

\eexam

The category of sheaves, $\sh(X)$, is an Abelian category contained in $\Psh(X)$; but it is not an Abelian subcategory: cokernels in the two categories differ.

For a topological space $X$, let $\O$ be the sheaf which maps $U$ to $\O(U)$ the Abelian group of continuous maps $U\to\bC$.
Similarly define $\O^\times(U)$ to be the Abelian group of continuous maps $U\to\bC^\times$.
Then we have a short exact sequence of sheaves:
$$ 0 \to \bZ \xto{2\pi i} \O \xto{\ \exp\ } \O^\times \to 0 $$
This is not exact in $\Psh(X)$, though, when $X=\bC^\times$.

\bdefn

    An additive functor $F\colon\A\to\B$ between Abelian categories is {\emphcolor left exact} (resp.\ {\emphcolor right exact}) if for every short exact sequence
    $0\to A\to B\to C\to0$ in $\A$, $0\to FA\to FB\to FC$ (resp.\ $FA\to FB\to FC\to0$) is exact in $\B$.
    $F$ is then {\emphcolor exact} if it is both left- and right-exact (meaning it preserves short exact sequences).

\edefn

\bexam

    The inclusion $\sh(X)\to\Psh(X)$ is left-exact (but not right-exact).
    There is an exact functor $\Psh(X)\to\sh(X)$ called {\emphcolor sheafification}.

\eexam

\bprop

    An additive functor $F\colon\A\to\B$ is left exact (resp.\ right exact) iff the exactness of $0\to A\to B\to C$ (resp.\ $A\to B\to C\to0$) implies the exactness
    of $0\to FA\to FB\to FC$ (resp.\ $FA\to FB\to FC\to0$).

\eprop

\bproof

    We prove this in categories of modules, and use Freyd-Mitchell to justify the extension to Abelian categories.
    $0\to A\to B\to C$ is exact iff $0\to A\to B\to C'\to0$ is exact where $C'$ is the image of the final map in $C$.
    If $F$ is left exact then $0\to FA\to FB\to FC'$ is exact, now we claim $FC'\to FC$ is a monomorphism and thus $0\to FA\to FB\to FC$ is exact as well.
    Indeed, $0\to C'\to C\to C/C'\to 0$ is exact and thus since $F$ is left exact $0\to FC'\to FC$ is exact, meaning $FC'\to FC$ is a monomorphism.

    Now suppose $F$ preserves exactness of sequences of the form $0\to A\to B\to C$, then in particular when we can append $\to0$.
    \qed

\eproof

\bprop

    Let $\A$ be an Abelian category, then $\hom_\A(M,\bul)$ is a left exact functor $\A\to\Ab$ for every $M\in\A$.

\eprop

\bproof

    Suppose we have $0\to A\xto fB\xto gC\to0$ in $\A$, we want to show that the following is exact
    $$ 0 \to \hom(M,A) \xto{\ f_*\ } \hom(M,B) \xto{\ g_*\ } \hom(M,C) $$
    where $f_*\alpha=f\circ\alpha$.
    Note that if $f_*\alpha=0$ then $\alpha=0$ since $f$ is monic, thus $f_*$ is monic.
    Furthermore since $g\circ f=0$, $g_*f_*(\alpha)=g\circ f\circ\alpha=0$ so $g_*f_*=0$.
    Now if $\beta\in\hom(M,B)$ is in $g_*$'s kernel then $g\circ\beta=0$ so $\beta(M)\subseteq\ker g=\im f$, thus $\beta$ factors through $f$, meaning it is in $f_*$'s image as desired.
    \qed

\eproof

\bcoro

    $\hom_\A(\bul,M)$ is left exact.

\ecoro

\bproof

    $\hom_\A(\bul,M)=\hom_{\A^\op}(M,\bul)$.

\eproof

\bdefn

    We define the {\emphcolor Yoneda embedding} $\y\colon\A\to[\A^\op,\Ab]$ for any additive category by mapping $A\mapsto h_A=\hom_\A(\bul,A)$ and $f\colon A\to B$ to $f_*\alpha=f\circ\alpha$.

\edefn

Considering the Yoneda embedding in classical category theory shows that this is a fully faithful embedding.
When $\A$ is Abelian $\y$ is left exact since $\y(A)=h_A$ is.

\blemm[title=Yoneda]

    The Yoneda embedding reflects exactness: $A\xto\alpha B\xto\beta C$ in $\A$ is exact if for every $M$ the following sequence is exact
    $$ \hom(M,A) \xto{\alpha_*} \hom(M,B) \xto{\beta_*} \hom(M,C) $$

\elemm

\bproof

    Let $M=A$ then we have $\beta\alpha=\beta_*\alpha_*\id=0$.
    Taking $M=\ker\beta$ we see that the inclusion $\iota\colon\ker\beta\to B$ satisfies $\beta_*\iota=\beta\iota=0$.
    Thus there is $\sigma\in\hom(M,A)$ such that $\iota=\alpha_*\sigma=\alpha\sigma$, so $\ker\beta=\im\iota\subseteq\im\alpha$ as desired.
    \qed

\eproof

